\par\Needspace{6\baselineskip}\medskip
\begin{center}\hypertarget{native-186-article139}{}\pdfbookmark[2]{139.}{article-139}139.\end{center}
Sint in sphaera coelesti $C,C',C''$ tria corporis coelestis loca heliocentrica,
quae resp. in circulis maximis $AB,A'B',A''B''$, et quidem inter $A$ et $B$, $A'$ et
$B'$, $A''$ et $B''$ sita erunt (art. 64, III): praeterea puncta $C,C',C''$ in eodem circulo
maximo iacebunt, puta in eo, quem planum orbitae in sphaera coelesti
proiicit. Designabimus per $r,r',r''$ tres corporis coelestis distantias a Sole;
per $\rho,\rho',\rho''$ eiusdem distantias a terra; per $R,R',R''$ terrae distantias a Sole.
Porro statuimus arcus $C'C'',CC'',CC'$ resp. $=2f,2f',2f''$, atque
\[
r'r''\sin2f=n,\qquad rr''\sin2f'=n',\qquad rr'\sin2f''=n''.
\]
\hypertarget{native-187}{}Habemus itaque
\[
\begin{gathered}
f'=f+f''\\[.4ex]
AC+CB=\delta,\qquad A'C'+C'B'=\delta',\\[.4ex]
A''C''+C''B''=\delta'',
\end{gathered}
\]
nec non
\[
\begin{aligned}
\frac{\sin\delta}{r}&=\frac{\sin AC}{\rho}=\frac{\sin CB}{R}\\[.7ex]
\frac{\sin\delta'}{r'}&=\frac{\sin A'C'}{\rho'}=\frac{\sin C'B'}{R'}\\[.7ex]
\frac{\sin\delta''}{r''}&=\frac{\sin A''C''}{\rho''}=\frac{\sin C''B''}{R''}\mathpunct{\cdot}
\end{aligned}
\]

Hinc patet, simulac situs punctorum $C,C',C''$ innotuerit, quantitates $r,r',r'',\rho,\rho',\rho''$
determinabiles fore. Iam ostendemus, quomodo ille e quantitatibus
\[
\frac{n''}{n}=P,\qquad 2\left(\frac{n+n''}{n'}-1\right)r'^3=Q
\]
elici possit, a quibus methodum nostram proficisci iam supra declaravimus.
